% Figure from MATH 452 notes, section 15. Compile with the notes preamble.
\begin{tikzpicture}[scale=1.0]
% ======== TOP LEFT: (r,theta) rectangle grid ========
\begin{scope}
\draw[->] (-0.2,0) -- (2.9,0) node[right, font=\small] {$r$};
\draw[->] (0,-0.2) -- (0,2.1) node[above, font=\small] {$\theta$};
\draw[thick, figB!70, fill=figB!6] (0.5,0.25) rectangle (2.45,1.72);
% grid: r lines at 0.5,0.89,1.28,1.67,2.06,2.45 ; theta at 0.25,0.615,0.98,1.355,1.72
\foreach \r in {0.89,1.28,1.67,2.06} \draw[figB!40, thin] (\r,0.25) -- (\r,1.72);
\foreach \t in {0.615,0.98,1.355} \draw[figB!40, thin] (0.5,\t) -- (2.45,\t);
% two same-size cells: inner-r and outer-r
\draw[thick, figG, fill=figG!22] (0.89,0.615) rectangle (1.28,0.98);
\draw[thick, figR, fill=figR!22] (2.06,0.615) rectangle (2.45,0.98);
\node[font=\scriptsize, align=center] at (1.45,-0.55) {equal cells in the $(r,\theta)$ plane};
\end{scope}
% map arrow
\draw[->, thick] (3.2,1.0) -- (4.6,1.0);
\node[font=\small, align=center] at (3.72,1.55) {$\Phi(r,\theta)=$\\ $(r\cos\theta,\, r\sin\theta)$};
% ======== TOP RIGHT: polar image ========
\begin{scope}[shift={(5.0,-0.1)}]
\draw[->] (-0.2,0) -- (3.3,0) node[right, font=\small] {$x$};
\draw[->] (0,-0.2) -- (0,2.7) node[above, font=\small] {$y$};
% sector between theta=0.25 and 1.72 rad, r from 0.5 to 2.45
% arcs at r values
\foreach \r in {0.5,0.89,1.28,1.67,2.06,2.45}
  \draw[figB!40, thin] ({\r*cos(14.3)},{\r*sin(14.3)}) arc[start angle=14.3, end angle=98.5, radius=\r];
% rays at theta values (deg): 14.3, 35.2, 56.2, 77.6, 98.5
\foreach \t in {14.3,35.2,56.2,77.6,98.5}
  \draw[figB!40, thin] ({0.5*cos(\t)},{0.5*sin(\t)}) -- ({2.45*cos(\t)},{2.45*sin(\t)});
% boundary of the sector (thick)
\draw[thick, figB!70]
  ({0.5*cos(14.3)},{0.5*sin(14.3)}) arc[start angle=14.3, end angle=98.5, radius=0.5]
  -- ({2.45*cos(98.5)},{2.45*sin(98.5)})
  arc[start angle=98.5, end angle=14.3, radius=2.45] -- cycle;
% image cells: inner (green) r in [0.89,1.28], theta in [35.2,56.2]; outer (red) r in [2.06,2.45]
\draw[thick, figG, fill=figG!22]
  ({0.89*cos(35.2)},{0.89*sin(35.2)}) arc[start angle=35.2, end angle=56.2, radius=0.89]
  -- ({1.28*cos(56.2)},{1.28*sin(56.2)})
  arc[start angle=56.2, end angle=35.2, radius=1.28] -- cycle;
\draw[thick, figR, fill=figR!22]
  ({2.06*cos(35.2)},{2.06*sin(35.2)}) arc[start angle=35.2, end angle=56.2, radius=2.06]
  -- ({2.45*cos(56.2)},{2.45*sin(56.2)})
  arc[start angle=56.2, end angle=35.2, radius=2.45] -- cycle;
\node[font=\scriptsize, align=right, anchor=east] at (0.42,-1.42) {images: the outer cell is wider ---\\ arc length $= r\,d\theta$ grows with $r$};
\end{scope}
% ======== BOTTOM: zoom of the outer cell ========
\draw[dashed, gray] (6.146,1.612) -- (5.09,-3.365);
\draw[dashed, gray] (7.002,1.312) -- (7.30,-3.22);
\begin{scope}[shift={(5.55,-4.4)}]
% magnified cell: nearly a rectangle, radial side dr, tangential side r0 dtheta
% at theta0=45.7deg: radial dir (0.698,0.716); tangential (-0.716,0.698)
\coordinate (O) at (0.6,0);
\coordinate (Ar) at (1.15,1.18);    % Phi_r dr direction scaled
\coordinate (At) at (-1.18,1.15);   % Phi_theta dtheta (longer: r0 dtheta)
\draw[thick, figR, fill=figR!14] (O) -- ($(O)+(Ar)$) -- ($(O)+(Ar)+0.9*(At)$) -- ($(O)+0.9*(At)$) -- cycle;
\draw[->, thick, figG] (O) -- ($(O)+(Ar)$) node[right, font=\small, yshift=-2pt] {$\boldsymbol{\Phi}_r\,dr$};
\draw[->, thick, figG] (O) -- ($(O)+0.9*(At)$) node[above left=-2pt, font=\small] {$\boldsymbol{\Phi}_\theta\,d\theta$};
% right angle mark
\draw[thin] ($(O)+0.16*(Ar)$) -- ($(O)+0.16*(Ar)+0.144*(At)$) -- ($(O)+0.144*(At)$);
\node[font=\small, figG!75!black] at (0.6,-0.62) {$\boldsymbol{\Phi}_r\,dr = (\cos\theta_0,\sin\theta_0)\,dr \;\perp\; \boldsymbol{\Phi}_\theta\,d\theta$};
\node[font=\small, figR!85!black] at (0.6,-1.08) {orthogonal sides $dr$ and $r_0\,d\theta$: \; area $= r_0\,dr\,d\theta$, \; so \; $|J| = r_0$};
\node[font=\scriptsize, gray] at (0.6,-1.50) {zoom: the cell straightens into a rectangle};
\end{scope}
\end{tikzpicture}
